\documentclass[11pt]{article}

\usepackage[margin=1in]{geometry}
\usepackage{amsmath}
\usepackage{booktabs}
\usepackage{graphicx}
\usepackage{hyperref}
\hypersetup{hidelinks}
% Edit these fields for your report.
\newcommand{\ReportTitle}{Statistics Analysis Report: [Study Title]}
\newcommand{\StudentName}{[Your Name]}
\newcommand{\CourseName}{[Course]}
\newcommand{\InstitutionName}{[Institution]}
\newcommand{\SubmissionDate}{[Submission Date]}

% Common statistics notation. Use each command in math mode.
\newcommand{\Prob}[1]{\operatorname{P}\!\left(#1\right)}
\newcommand{\Expect}[1]{\operatorname{E}\!\left[#1\right]}
\newcommand{\Variance}[1]{\operatorname{Var}\!\left(#1\right)}
\newcommand{\SampleMean}[1]{\overline{#1}}


\title{\ReportTitle}
\author{\StudentName \\ \CourseName \\ \InstitutionName}
\date{\SubmissionDate}

\begin{document}
\maketitle

\begin{abstract}
State the research question, data source, method, and main finding in a few sentences. The values in this sample are synthetic and demonstrate formatting only.
\end{abstract}

\section{Introduction}
Describe the question you are investigating and why it matters. Define the outcome and the groups or variables you will compare.

\section{Data and methods}
Describe where your real data came from, how observations were selected, and any cleaning or exclusions. Name the summary statistics or model you use. For example, the sample mean is
\begin{equation}
  \SampleMean{x} = \frac{1}{n}\sum_{i=1}^{n} x_i.
  \label{eq:mean}
\end{equation}
For a random variable $X$, use $\Expect{X}$ and $\Variance{X}$ when reporting population quantities. Explain why your chosen method addresses the research question and state any relevant assumptions.

\section{Results}
Table~\ref{tab:summary} shows an illustrative layout. Every number in the table is synthetic; replace the caption, labels, and values with your results.

\begin{table}[htbp]
  \centering
  \caption{Synthetic descriptive statistics for two example groups}
  \label{tab:summary}
  \begin{tabular}{lrrr}
    \toprule
    Group & Sample size & Mean & Standard deviation \\
    \midrule
    A & 10 & 12.4 & 2.1 \\
    B & 10 & 15.1 & 2.4 \\
    \bottomrule
  \end{tabular}
\end{table}

Figure~\ref{fig:placeholder} marks where a plot can go. Replace the box with an image using \verb|\includegraphics| and give the plot an informative caption.

\begin{figure}[htbp]
  \centering
  \fbox{\parbox[c][1.5in][c]{0.7\textwidth}{\centering Replace this box with a plot of your results.}}
  \caption{Placeholder for a figure; no data are plotted here}
  \label{fig:placeholder}
\end{figure}

\section{Discussion}
Interpret your results in context. Address uncertainty, assumptions, limitations, and whether the data support the question posed in the introduction. Do not infer a real-world effect from the synthetic example values.

\section{Conclusion}
Summarize the answer to your research question and identify a useful next step.

\section*{References}
[Add the sources you actually used in the format required by your course.]

\end{document}
